%%
% 第六章 总结与展望（~2页）
%%

\begin{conclusion}

\section{工作总结}

本文针对学术多模态属性图数据集的空白，提出了AMMAG——首个面向学术论文领域的多模态属性图Benchmark数据集及其基于多智能体的自动化构建系统。本文的主要工作和贡献总结如下：

\textbf{（1）构建了首个学术多模态属性图Benchmark数据集AMMAG。}AMMAG覆盖计算机科学、数学和物理学三个学科大类，包含cs.SI（社交信息学）、cs.LO（计算逻辑）、math.LO（数理逻辑）和physics.soc-ph（社会物理）四个子领域，共计45,672篇论文。每篇论文节点包含全文文本（Markdown格式）、实验图表（PNG格式）和数学公式（LaTeX格式）三种模态属性，并通过59,364条引用边构成引用网络图结构。与现有MAG基准相比，AMMAG首次将公式模态引入MAG学习范式；与OGB\cite{hu2020ogb}和CS-TAG\cite{yan2023cstag}等学术图基准相比，AMMAG提供了原始多模态属性而非低维压缩向量表示。四个子领域在模态分布上的显著差异（如math.LO的公式密度是cs.SI的6.3倍）为模态消融实验提供了天然的实验条件。

\textbf{（2）设计了基于多智能体的自动化构建系统。}本文设计并实现了包含五个专用智能体的自动化构建系统，采用编排器-智能体-工具三层架构。五个智能体分别负责论文下载、多模态提取、数据清洗、引用图构建和标准化格式化，实现了从arXiv元数据到标准化MAG数据集的端到端自动化流程。其中，多模态提取引擎基于DocLayout-YOLO\cite{zhao2024doclayoutyolo}、UniMERNet\cite{wang2024unimernet}和PaddleOCR\cite{du2020ppocr}实现了文本、图像和公式的并行提取；引用图构建模块通过双源引用提取（PDF+Markdown，209\%提升）和三级匹配引擎（arXiv ID直接匹配、精确标题匹配、倒排索引模糊匹配）实现了自动化引用网络构建。并行化策略通过多GPU批量提取、多进程并行清洗和fork-based COW零拷贝优化实现了高效的数据处理。

\textbf{（3）进行了全面的Benchmark实验评估。}本文在多模态消融、模型有效性、鲁棒性、Few-shot学习、链接预测和跨域迁移六组任务上，系统评估了13种模型（含7种标准GNN/MLP和6种多模态GNN）在7种模态组合下的性能表现，采用CLIP-ViT-Large统一768维编码方案，共计产生超过1,000组实验结果。实验结果表明，GCNII以91.09\%取得最优节点分类准确率，但Top-tier模型间差异不显著（$p > 0.05$），约91\%的性能天花板需要从特征工程层面突破。

\textbf{（4）发现了学术领域多模态特性的重要规律。}本文首次在学术图数据集上量化分析了文本、图像和公式三种模态的贡献差异。核心发现包括：（a）T+F（文本+公式）是最佳模态组合（均值90.84\%），甚至略优于三模态T+I+F（90.62\%），公式比图像提供了更多互补增益——这是MAGB\cite{yan2025magb}等现有基准因缺少公式模态无法得出的独有发现；（b）不同学术领域对模态的偏好差异显著，cs.SI从多模态获益最大（+2.11pp）而math.LO从公式获益最大（+1.43pp）；（c）跨域迁移增益接近零（$\pm$0.5pp），领域间子分类模式差异大，领域特定训练是必要的；（d）链接预测在该数据集上饱和（MRR $>$ 0.95），不适合作为主要评估指标；（e）CLIP统一编码方案远优于独立编码器（跨模态检索MRR提升67倍）。

\section{不足与展望}

尽管本文在学术多模态属性图数据集的构建和评估方面取得了一定进展，但仍存在以下不足，每条不足对应一个未来研究方向：

\textbf{（1）数据源单一。}当前AMMAG仅覆盖arXiv单一数据源，以计算机科学、数学和物理学领域的论文为主，未涵盖生物医学、社会科学等重要学科。未来可将系统扩展至PubMed\cite{lo2020s2orc}（生物医学）、DBLP（计算机科学完整覆盖）、Semantic Scholar（跨学科）等数据源，进一步增加数据集的学科覆盖面和规模。不同数据源的论文格式和元数据标准存在差异，这需要对系统的提取和匹配模块进行适配。

\textbf{（2）引用边匹配率有待提升。}当前引用边匹配率约20.7\%，大量跨领域引用因不在候选集内而丢失。三级匹配引擎在处理OCR错误和非标准引用格式时仍存在误匹配的情况。未来可引入大语言模型（LLM）辅助引用匹配——利用LLM的语义理解能力对模糊引用进行消歧，提高匹配的准确率和召回率。此外，可以结合Semantic Scholar\cite{lo2020s2orc}和OpenAlex\cite{priem2022openalex}等外部学术知识库的引用关系数据，补充通过PDF解析未匹配到的引用边。

\textbf{（3）多模态融合策略较初步。}当前采用拼接+可学习线性投影作为多模态融合策略，实验表明T+F已接近T+I+F的性能（90.84\% vs 90.62\%），说明图像模态的互补信息有限。未来可探索更先进的融合方法：注意力融合机制（Cross-Modal Attention）能够自适应地学习不同模态的权重；多模态图Transformer可以在消息传递过程中建模模态间的交互关系；基于视觉语言模型（如LLaVA、Qwen-VL等）的统一编码器可以从根本上消除不同模态间的表示差异。

\textbf{（4）数据集规模相对较小。}AMMAG当前包含45,672篇论文，相比OGB的ogbn-arxiv（169K节点）规模较小。规模限制可能影响大规模图学习方法的评估效果。未来可通过以下途径扩展数据集规模：增加每个领域的论文数量（如从12,000篇扩展至50,000篇）、增加子领域的覆盖范围（如增加cs.AI、cs.CV等热门CS子领域）、以及引入跨arXiv的论文数据。在扩展规模的同时，需要保证多模态提取的质量——大规模数据处理对系统的并行能力和容错机制提出了更高要求。

\textbf{（5）图像编码器不适合科学图表。}虽然本文使用了CLIP统一编码方案（768维），在跨模态对齐上远优于独立编码器（检索MRR提升67倍），但CLIP视觉编码器在自然图像上预训练，其视觉特征不适用于科学图表（架构图、实验结果图、流程图等），Image-only准确率仅约81\%，远低于Text-only的约89\%。未来可训练领域适配的科学图表编码器：利用DocGenome\cite{xia2024docgenome}等大规模科学文档数据集微调视觉模型，使编码器学习科学图表的特有视觉特征；或使用多模态科学模型进行端到端的图表特征提取，将图表的视觉信息与上下文文本信息联合编码。

\textbf{（6）公式编码器仍有提升空间。}本文将LaTeX公式作为文本序列用CLIP文本编码器编码，虽已显著优于SBERT编码方案，但公式模态的单模态性能（约78\%）仍远低于文本（约89\%）。Few-shot实验发现公式特征需要更多样本释放潜力（k=1$\to$20提升26.4pp，相对提升70.8\%）。未来应引入MathBERT等数学感知编码器进行对比实验，基于公式树（AST）的编码方法可能更有效地捕获公式的结构化语义。

\end{conclusion}
